\documentclass[10pt,a4paper]{amsart}
\usepackage[margin=19mm]{geometry}
\usepackage{amsmath,amssymb,mathtools}
\usepackage[scr=rsfs,cal=euler]{mathalfa}
\usepackage{xfrac}
\usepackage[unicode,hidelinks,bookmarks=false]{hyperref}
\newcommand{\ie}{i.e.\ }
\newcommand{\cf}{cf.\ }
\newcommand{\resp}{respectively\ }
\newcommand{\Subsection}{\subsection}
\providecommand{\nobreakdash}{\nobreak\hbox{-}}
\setlength{\emergencystretch}{2em}
\multlinegap0pt
\usepackage{bm}
\usepackage{bbm}
\usepackage{dsfont}
\usepackage{xspace}
\usepackage{cancel}
\usepackage{mathtools}
\usepackage{amsthm}
\makeatletter
\@ifundefined{thm}{\newtheorem{thm}{Thm}}{}
\@ifundefined{prop}{\newtheorem{prop}[thm]{Proposition}}{}
\makeatother
\DeclareMathOperator{\re}{Re}
\newcommand{\C}{\ensuremath{\mathbb{C}}}
\title[Spectral Moment Formulae for $GL(3)\times GL(2)$ $L$-functio]{Spectral Moment Formulae for $GL(3)\times GL(2)$ $L$-functions III: The Twisted Case}
\author{Chung-Hang Kwan}
\date{Source: arXiv:2311.15417v2 (CC BY 4.0)}
\begin{document}
\maketitle

\noindent\emph{Source and licence.} Spectral Moment Formulae for $GL(3)\times GL(2)$ $L$-functions III: The Twisted Case. Version arXiv:2311.15417v2 (CC BY 4.0), distributed under the Creative Commons Attribution 4.0 International licence, as recorded in the arXiv licence metadata for this exact version and shown on the abstract page https://arxiv.org/abs/2311.15417v2. Full rights notice: licenses/arxiv-2311.15417-CC-BY-4.0.txt.\par\smallskip

\noindent\textit{Editorial scope.} Counted unit: the Prop whose source label is \texttt{cuspdual} in the article (source0.tex lines 1662--1673, source label \texttt{cuspdual}), delivered together with the article's own complete original proof of it (source0.tex lines 1676--1734, one \texttt{proof} environment of 59 lines). No mathematics is written, repaired or re-derived: statement and proof are the article's own text line for line. Required context retained uncounted: regAC (lines 1450--1452); congHurw (lines 1486--1488); twistDS (lines 1494--1496); gl3EPid (lines 1603--1605). Every definition, display and label of those lines is retained unchanged. Omitted from this package and not redistributed: 1--1449, 1453--1485, 1489--1493, 1497--1602, 1606--1661, 1674--1675, 1735--2211. Nothing inside a retained excerpt is omitted or altered: every retained body is delivered exactly as the article's lines are written, so no omission receipt of any kind accompanies this package. Citations: the retained excerpts carry no citation command at all, so no bibliography entry of the article is transcribed and no reference is redistributed. Reference adaptation: none was needed and none was made; every reference inside the delivered unit resolves inside it, so no unresolved reference or citation is produced. Printed slips kept literal: no printed word, symbol, hypothesis or number of the counted statement or of its proof was altered. Layout and class: the article's own class is \texttt{amsart} with packages including hyperref, amsmath, amsthm, amssymb, color, url, soul, empheq, mathdots, bm, mathabx, epsfig, pst-node, tikz-cd; the wrapper is the lane's pinned \texttt{amsart} editorial wrapper at 10pt on A4, which restates the theorem-like environments the retained text needs and adds the article's own macro definitions that the retained text needs (\texttt{\textbackslash C}, \texttt{\textbackslash re}), with extra wrapper packages (bm, bbm, dsfont, xspace, cancel, mathtools). The delivered numbering is the unit's own. Assets: no retained span contains a figure environment, a graphics-inclusion command, a PDF-page inclusion, a table or a TikZ picture; no image file is copied into this package, the package declares no asset dependency and no image file is redistributed with it. Licence: version arXiv:2311.15417v2 of this article is distributed under the Creative Commons Attribution 4.0 International licence, as recorded in the arXiv licence metadata for this exact version and shown on the abstract page https://arxiv.org/abs/2311.15417v2. A full rights notice is delivered at licenses/arxiv-2311.15417-CC-BY-4.0.txt. Characters: every retained excerpt is pure ASCII, so the delivered engine, XeLaTeX with its default Latin Modern text font, reports no missing character.
\begin{align}\label{regAC}
	\re (2s-s_{0}) \ > \ 1 \hspace{20pt} \text{ and  } \hspace{20pt}\re s_{0} \ > \ 1+\theta. 
\end{align}

\begin{align}\label{congHurw}
	\sum_{\substack{a_{0} \ge 1 \\ a_{0} \ \equiv \  \ell \ (\bmod\, d)}}  \frac{1}{a_{0}^{2s-s_{0}}} \ &= \  d^{-(2s-s_{0})}\cdot  \zeta\left( 2s-s_{0}, \frac{\ell}{d}\right),
\end{align}

 \begin{align}\label{twistDS}
 	\mathcal{L}^{(a)}_{\pm}\left(s_{0}, s; \Phi \right)  \ \ = \ \    \sum_{dr \mid a} \   \frac{d^{2s_{0}-1} \mu(r)}{ r^{s_{0}} }\ \overline{\lambda_{\Phi}}\left( \frac{a}{dr}\right)  \  \sideset{}{^*}{\sum}_{\ell \ (\bmod\, d)} \  \zeta\left( 2s-s_{0}, \frac{\ell}{d}\right) 	L\left( s_{0}; \mp \frac{r\overline{\ell}}{d}; \Phi\right)
 \end{align}

  \begin{align}\label{gl3EPid}
  	 L(s, \Phi) \ &= \  \prod_{p} \ \left(1-\lambda_{\Phi}(p)p^{-s} +\overline{\lambda_{\Phi}}(p)p^{-2s}-p^{-3s}\right)^{-1} \hspace{20pt} (\re s \ \gg \ 1)
  \end{align}

\begin{prop}\label{cuspdual}
	For any $(s_{0},s)\in \C^2$ except on the line $2s-s_{0}=1$, we have
	\begin{align}\label{expinmul}
		\mathcal{L}^{(p)}_{\pm}\left(s_{0}, s; \Phi \right)  \  \ =  \ \ 	& 	\mathcal{P}_{p}(s_{0}, s; \Phi) \,  \zeta(2s-s_{0}) \, L(s_{0}, \Phi) \nonumber\\
		& \hspace{20pt} \ + \  	\frac{ p^{ s_{0}+2s-1}}{\phi(p)}  \    \left( \ \  \sideset{}{^+}{\sum}_{\chi \ (\bmod\, p)} \ \  \mp \ \  \sideset{}{^-}{\sum}_{\chi \ (\bmod\, p)} \right) \ \tau(\chi) L(2s-s_{0}, \chi)  L\left(s_{0}, \Phi\otimes \overline{\chi}\right),
	\end{align}
	where   $\tau(\chi):= \sideset{}{^*}{\sum}\limits_{x \, (\bmod\, p)} \chi (x) e(x/p)$ is the Gauss sum of\, $\chi \, (\bmod\, p)$ and 
	\begin{align}
		\mathcal{P}_{p}(s_{0}, s; \Phi) \ \ := \ \  &\frac{p^{2s+s_{0}-1}}{\phi(p)} \, (1-p^{-(2s-s_{0})}) \left(-1 +\lambda_{\Phi}(p)p^{1-s_{0}}- \overline{\lambda_{\Phi}}(p) p^{1-2s_{0}}+p^{1-3s_{0}}\right)\nonumber\\
		 &\hspace{40pt}  \ + \ \overline{\lambda_{\Phi}}(p) \ -  \  p^{-s_{0}}.  
	\end{align}
\end{prop}

\begin{proof}
	The computations below will be first  carried out on the region of absolute convergence (\ref{regAC}). Then  our desired conclusion will follow from analytic continuation.  In  (\ref{twistDS})  we either have $d=1$ or $d=p$,  and  correspondingly   $r\mid p$ or $r\mid 1$. 
	
	
	\noindent \textbf{Case (1): $d=1$. }  We take $\ell=1$ in this case and then 
	\begin{align*}
		\mathcal{L}^{(p)}_{\pm}\left(s_{0}, s; \Phi \right) \bigg|_{d=1} \ \ = \ \  	&\ \sum_{r\in \{1,p\}} \ \frac{\mu(r)}{r^{s_{0}}}\cdot   \overline{\lambda_{\Phi}}(p/r) \cdot \zeta(2s-s_{0}) L\left(s_{0}, \mp r; \Phi\right) \nonumber\\
		\ \ =  \ \ &  \left( \,  \overline{\lambda_{\Phi}}(p)- p^{-s_{0}}\right) \zeta(2s-s_{0}) L\left(s_{0}, \Phi\right).
	\end{align*}
	
	
	\noindent \textbf{Case (2): $d=p$.} We have
	\begin{align}\label{case2}
		\mathcal{L}^{(p)}_{\pm}\left(s_{0}, s; \Phi\right) \bigg|_{d=p}  \ = \ 	p^{ 2s_{0}-1}  \  \   \sideset{}{^*}{\sum}_{\ell \ (\bmod\, p)} \  \zeta\left( 2s-s_{0}, \frac{\ell}{p}\right) 	L\left( s_{0}; \mp \frac{\overline{\ell}}{p}; \Phi\right). 
	\end{align}
The Hurwitz zeta function in (\ref{case2})   can be rewritten as  
	\begin{align*}
		\zeta\left( 2s-s_{0}, \frac{\ell}{p}\right) \  \ &=  \ \   \frac{p^{2s-s_{0}}}{\phi(p)} \ \sum_{\chi \ (\bmod\, p)} \ \overline{\chi}(\ell) L(2s-s_{0}, \chi)
	\end{align*}
	 with (\ref{congHurw}) and   the orthogonality relation for Dirichlet characters.  Hence, we have
	\begin{align}\label{afterorth}
			\mathcal{L}^{(p)}_{\pm}\left(s_{0}, s; \Phi\right) \bigg|_{d=p}  \ &= \ 	 \frac{p^{2s+s_{0}-1}}{\phi(p)}  \      \sum_{\chi \ (\bmod\, p)} \  L(2s-s_{0}, \chi)  \  \sideset{}{^*}{\sum}_{\ell \ (\bmod\, p)}  \ \overline{\chi}(\ell) L\left( s_{0}; \mp \frac{\overline{\ell}}{p}; \Phi\right) \nonumber\\
			\ &=\   	 \frac{p^{2s+s_{0}-1}}{\phi(p)}  \      \sum_{\chi \ (\bmod\, p)} \  L(2s-s_{0}, \chi)  \ \sum_{n=1}^{\infty} \ \frac{\lambda_{\Phi}(n)}{n^{s_{0}}} \  \sideset{}{^*}{\sum}_{\ell \ (\bmod\, p)}  \ \overline{\chi}(\ell) 	e\left(\mp \frac{n\bar{\ell}}{p}\right).
	\end{align}
	
	If $(n,p)=1$, then the replacement\,   $\ell \to \mp n\bar{\ell}$\, $(\bmod\, p)$ implies 
	\begin{align}\label{intogauss}
		\sideset{}{^*}{\sum}_{\ell \ (\bmod\, p)}  \ \overline{\chi}(\ell) 	e\left(\mp \frac{n\bar{\ell}}{p}\right) \ \ =  \ \  \overline{\chi}(\mp n) \tau(\chi). 
	\end{align}
 (When $\chi=\chi_{0}$, equation (\ref{intogauss}) is equal to $-1$.) If  $p\mid n$, then
	\begin{align*}
		\sideset{}{^*}{\sum}_{\ell \ (\bmod\, p)}  \ \overline{\chi}(\ell) 	e\left(\mp \frac{n\bar{\ell}}{p}\right) \ \ = \  \  	\sideset{}{^*}{\sum}_{\ell \ (\bmod\, p)}  \ \overline{\chi}(\ell)  \ \ = \ \  \phi(p)\cdot  1_{\chi=\chi_{0}}. 
	\end{align*}
	
	As a result, the contribution from $\chi=\chi_{0}$ in (\ref{afterorth}) is equal to 
	\begin{align}
			\mathcal{L}^{(p)}_{\pm}\left(s_{0}, s; \Phi\right) \bigg|_{d=p, \, \chi=\chi_{0}}   \ \ &= \  \   \frac{p^{2s+s_{0}-1}}{\phi(p)} \,  L(2s-s_{0}, \chi_{0})  \left\{ -\sum_{(n,p)=1} \ \frac{\lambda_{\Phi}(n)}{n^{s_{0}}}  \ + \ \phi(p) \, \sum_{p\mid n} \ \frac{\lambda_{\Phi}(n)}{n^{s_{0}}}   \right\} \nonumber\\
		\ \ &= \ \     \frac{p^{2s+s_{0}-1}}{\phi(p)} \,  L(2s-s_{0}, \chi_{0})  \left\{  \phi(p) L(s_{0}, \Phi) \ - \  p \sum_{(n,p)=1} \ \frac{\lambda_{\Phi}(n)}{n^{s_{0}}} \right\}. \nonumber
	\end{align}
	It follows from (\ref{gl3EPid}) that
	\begin{align}
			\mathcal{L}^{(p)}_{\pm}\left(s_{0}, s; \Phi\right) \bigg|_{d=p, \, \chi=\chi_{0}} 	\ \ &= \  \     \frac{p^{2s+s_{0}-1}}{\phi(p)} \,  L(2s-s_{0}, \chi_{0})  \, L(s_{0}, \Phi) \nonumber\\
		& \hspace{60pt} \,  \cdot \,  \left\{  \phi(p) \ - \  p   \left(1-\lambda_{\Phi}(p)p^{-s_{0}} +\overline{\lambda_{\Phi}}(p)p^{-2s_{0}}-p^{-3s_{0}}\right) \right\}  \nonumber\\
		\ \ &= \  \  \frac{p^{2s+s_{0}-1}}{\phi(p)} \, (1-p^{-(2s-s_{0})}) \left(-1 +\lambda_{\Phi}(p)p^{1-s_{0}}- \overline{\lambda_{\Phi}}(p) p^{1-2s_{0}}+p^{1-3s_{0}}\right) \nonumber\\
		&\hspace{60pt} \cdot \,  \zeta(2s-s_{0}) \, L(s_{0}, \Phi). \nonumber
	\end{align}
	
	On the other hand,  the contributions from $\chi\neq \chi_{0}$ in (\ref{afterorth}) is equal to 
	\begin{align}
			\mathcal{L}^{(p)}_{\pm}\left(s_{0}, s; \Phi\right) \bigg|_{d=p, \, \chi\neq \chi_{0}}  \ \ &= \  \   \frac{p^{2s+s_{0}-1}}{\phi(p)}  \     \sideset{}{^*}{\sum}_{\chi \, (\bmod\, p)} \  \chi(\mp 1)\tau(\chi) L(2s-s_{0}, \chi)  \, \sum_{(n,p)=1} \ \frac{\lambda_{\Phi}(n)\overline{\chi}(n)}{n^{s_{0}}}   \nonumber\\
			\ \ &= \ \   \frac{p^{2s+s_{0}-1}}{\phi(p)}  \     \sideset{}{^*}{\sum}_{\chi \, (\bmod\, p)} \  \chi(\mp 1)\tau(\chi) L(2s-s_{0}, \chi) L(s_{0}, \Phi\otimes \overline{\chi}). \nonumber
	\end{align}
	The conclusion of Proposition \ref{cuspdual} follows readily from
	\begin{align*}
			\mathcal{L}^{(p)}_{\pm}\left(s_{0}, s; \Phi \right)  \ \ = \ \  	\mathcal{L}^{(p)}_{\pm}\left(s_{0}, s; \Phi \right) \bigg|_{d=1} \ + \  	\mathcal{L}^{(p)}_{\pm}\left(s_{0}, s; \Phi\right) \bigg|_{d=p, \, \chi=\chi_{0}}   \ + \ 	\mathcal{L}^{(p)}_{\pm}\left(s_{0}, s; \Phi\right) \bigg|_{d=p, \, \chi\neq \chi_{0}}  
	\end{align*}
	and   analytic continuation. 

\end{proof}

\end{document}
